\notepage{49}{Orbital eigenfunctions in position space}
The highest-weight state satisfies $L_+\ket{\ell,\ell}=0$.
To use this in position space, first express the operators in spherical
coordinates. With $\psi(\mathbf r)=\braket{\mathbf r}{\psi}$,
\[
 \mel{\mathbf r}{L_x}{\psi}
 =-\ii\hbar\left(y\frac{\partial}{\partial z}
                 -z\frac{\partial}{\partial y}\right)\psi(\mathbf r).
\]
\editorial{The Cartesian prefactor is $-\ii\hbar$, from
$\mathbf P=-\ii\hbar\boldsymbol\nabla$; the handwritten line has the
opposite sign.}
Using $y=r\sin\theta\sin\phi$, $z=r\cos\theta$,
\[
 \cos\phi=\frac{x}{\sqrt{x^2+y^2}},\qquad
 \cos\theta=\frac{z}{r},
\]
the chain rule gives
\begin{align*}
 y\partial_z-z\partial_y
 &=r\sin\theta\sin\phi
   \left(\cos\theta\partial_r-\frac{\sin\theta}{r}\partial_\theta\right)\\
 &\quad-r\cos\theta
   \left(\sin\theta\sin\phi\partial_r
    +\frac{\cos\theta\sin\phi}{r}\partial_\theta
    +\frac{\cos\phi}{r\sin\theta}\partial_\phi\right)\\
 &=-\sin\phi\partial_\theta-\cos\phi\cot\theta\partial_\phi.
\end{align*}
Thus
\begin{align*}
 L_x&=\ii\hbar\left(\sin\phi\partial_\theta
                         +\cos\phi\cot\theta\partial_\phi\right),\\
 L_y&=\ii\hbar\left(-\cos\phi\partial_\theta
                         +\sin\phi\cot\theta\partial_\phi\right),\\
 L_\pm&=\pm\hbar\ee^{\pm\ii\phi}
          \left(\partial_\theta\pm\ii\cot\theta\partial_\phi\right).
\end{align*}
The $L_y$ expression also follows by a $\pi/2$ azimuthal rotation of
the $L_x$ expression.

Since $L_z\ket{\ell,\ell}=\ell\hbar\ket{\ell,\ell}$, write
$\braket{r,\theta,\phi}{\ell,\ell}
=u_\ell^\ell(r,\theta)\ee^{\ii\ell\phi}$. Then
\[
 (\partial_\theta-\ell\cot\theta)u_\ell^\ell=0,
 \qquad
 \frac{\dd u_\ell^\ell}{u_\ell^\ell}
 =\ell\frac{\cos\theta}{\sin\theta}\dd\theta,
\]
so
\[
 \psi_\ell^\ell(r,\theta,\phi)=R_\ell(r)\sin^\ell\theta\,
 \ee^{\ii\ell\phi}.
\]
\textit{Exercise:} starting from $L_-\ket{\ell,-\ell}=0$, obtain
$\psi_\ell^{-\ell}=R_\ell(r)\sin^\ell\theta\ee^{-\ii\ell\phi}$,
up to a phase convention.
